\PassOptionsToPackage{dvipsnames,table}{xcolor}
\documentclass[11pt,a4paper]{article}

\usepackage{Act}
\begin{document}
\input{\detokenize{/home/fenarius/Travail/Cours/cpge-info/latex/Macros.tex}}
\ModeExercice
\setboolean{corrige}{true}
\IC{1}{Ligne de commande -- Premiers pas en C}
\setcounter{Exercise}{0}

\begin{Exercise}[title={QCM sur le système et le langage C}] \points{4}\\
	Dans cet exercice, chaque question peut avoir \textit{zéro, une ou plusieurs bonnes réponses}. Pour chaque question, cocher la ou les cases correspondantes aux réponses exactes.

	\Question{À propos du système de fichiers et des commandes Linux :}
	\begin{itemize}[itemsep=0pt,topsep=1pt,parsep=0pt]
		\item[\gc] Le chemin {\tt /home/alice/tp1} est un chemin absolu.
		\item[\bc] La commande {\tt cd ..} permet de revenir dans le répertoire personnel de l'utilisateur.
		\item[\gc] La commande {\tt ls -a} affiche tous les fichiers du dossier, y compris les fichiers cachés.
		\item[\bc] La commande {\tt rm} permet de renommer un dossier ou un fichier.
	\end{itemize}

	\Question{Concernant les liens sous Linux :}
	\begin{itemize}[itemsep=0pt,topsep=1pt,parsep=0pt]
		\item[\bc] Lorsqu'on crée un lien physique vers un fichier, on effectue une copie de ce fichier.
		\item[\bc] La commande {\tt ln source.txt lien} crée un lien symbolique.
		\item[\gc] Supprimer le fichier source rend un lien symbolique orphelin (inutilisable).
		\item[\bc] Deux fichiers ne peuvent jamais avoir le même \textit{inode}.
	\end{itemize}

	\Question{À propos des types de base en langage C :}
	\begin{itemize}[itemsep=0pt,topsep=1pt,parsep=0pt]
		\item[\gc] Le type {\tt int} permet de représenter des entiers relatifs.
		\item[\bc] Les chaînes de caractères possèdent le type natif {\tt str}.
		\item[\bc] L'opérateur {\tt **} permet d'élever un {\tt float} à une puissance.
		\item[\gc] Les types {\tt float} et {\tt double} permettent tous les deux de représenter des nombres décimaux.
	\end{itemize}

	\Question{À propos des types et de la syntaxe en langage C :}
	\begin{itemize}[itemsep=0pt,topsep=1pt,parsep=0pt]
		\item[\gc] L'inclusion de l'en-tête {\tt <stdbool.h>} est nécessaire pour utiliser le type {\tt bool}.
		\item[\gc] Une variable de type {\tt char} s'écrit entre apostrophes simples (ex : {\tt 'a'}), tandis qu'une chaîne de caractères s'écrit entre guillemets (ex : {\tt "a"}).
		\item[\bc] L'opérateur d'égalité en C s'écrit {\tt =} et l'affectation s'écrit {\tt :=}.
		\item[\bc] L'opérateur logique \og et \fg{} s'écrit {\tt ||} et l'opérateur logique \og ou \fg{} s'écrit {\tt \&\&}.
	\end{itemize}

	\Question{Pour se déplacer dans son répertoire personnel, l'utilisateur {\tt toto} peut taper :}
	\begin{itemize}[itemsep=0pt,topsep=1pt,parsep=0pt]
		\item[\bc] {\tt cd .}
		\item[\gc] {\tt cd \textasciitilde}
		\item[\bc] {\tt cd ..}
		\item[\gc] {\tt cd /home/toto/}
	\end{itemize}

	\Question{Que vaut la variable {\tt x} après l'instruction {\tt float x = 7 / 2;} en C ?}\\
	\begin{tabularx}{\linewidth}{@{}XXXX@{}}
		\bc {\tt 3.5} & \gc {\tt 3.0} & \bc {\tt 3} & \bc {\tt 1.0} \\
	\end{tabularx}

	\Question{Quel affichage produit le fragment de code suivant ?}
	\begin{langageC}
int s = 0;
for (int i = 1; i <= 4; i++) {
    s += i;
}
printf("%d\n", s);
	\end{langageC}
	\begin{tabularx}{\linewidth}{@{}XXXX@{}}
		\bc {\tt 4} & \bc {\tt 6} & \gc {\tt 10} & \bc {\tt 15} \\
	\end{tabularx}

	\Question{À propos des fonctions en langage C :}
	\begin{itemize}[itemsep=0pt,topsep=1pt,parsep=0pt]
		\item[\gc] En C, le passage des arguments scalaires se fait par valeur (une copie est transmise).
		\item[\bc] Une fonction ne peut jamais avoir {\tt void} comme type de retour.
		\item[\bc] Une fonction doit obligatoirement comporter au moins un argument.
		\item[\bc] Le mot-clé {\tt return} est obligatoire dans le corps de toute fonction.
	\end{itemize}
\end{Exercise}

\newpage
\begin{Exercise}[title={Ligne de commande et arborescence}] \points{6}\\
	Bob travaille dans un sous-dossier {\tt TP1} de son répertoire personnel. Ce dossier contient des fichiers source C ({\tt .c}) et des fichiers de données texte ({\tt .txt}). Écrire les commandes du shell qui lui permettent de réaliser les opérations suivantes :
	
	\Question{Afficher le chemin absolu du répertoire de travail actuel dans la console.}
	\reponse{0}{0.5}
	\tcor{\tt pwd}

	\Question{Créer un nouveau sous-dossier nommé {\tt src} dans le répertoire courant.}
	\reponse{0}{0.5}
	\tcor{\tt mkdir src}

	\Question{Déplacer tous les fichiers portant l'extension {\tt .c} du dossier courant dans le sous-dossier {\tt src}.}
	\reponse{0}{1}
	\tcor{\tt mv *.c src}

	\Question{Afficher les 10 premières lignes du fichier {\tt data.txt} situé dans le répertoire courant.}
	\reponse{0}{0.5}
	\tcor{\tt head -n 10 data.txt}


	\Question{Rediriger la liste détaillée des fichiers du répertoire courant (avec {\tt ls -l}) dans un fichier nommé {\tt contenu.txt}.}
	\reponse{0}{0.5}
	\tcor{\tt ls -l > contenu.txt}

	\Question{Créer un lien symbolique nommé {\tt raccourci.txt} dans son répertoire personnel ({\tt \textasciitilde}) qui pointe vers le fichier {\tt data.txt} du dossier courant.}
	\reponse{0}{1}
	\tcor{\tt ln -s \textasciitilde/TP1/data.txt \textasciitilde/raccourci.txt}

	\Question{En utilisant un tube (\textit{pipe}) {\tt |}, compter le nombre de fichiers et dossiers présents dans le répertoire courant (on rappelle que la commande {\tt wc -l} compte le nombre de lignes reçues en entrée).}
	\reponse{0}{1}
	\tcor{\tt ls | wc -l}

	\Question{Afficher la 42\textsuperscript{e} ligne du fichier {\tt prog.c} situé dans le répertoire courant.}
	\reponse{0}{1}
	\tcor{\tt head -n 42 prog.c | tail -n 1}


\end{Exercise}

\end{document}

